\section{系统族：宏观规律的压缩--闭合判据}
\label{sec:criterion}

上一节为每个有限系统给出了一条达到的复杂度--闭合误差前沿。现在考虑系统规模
增长时，这条前沿能否同时逼近低相对复杂度与低闭合误差。设
\[
(X_N,\cK_N,u_N,g_N),
\qquad |B_N|\longrightarrow\infty,
\]
并沿用第~\ref{sec:operational}节的可辨微观空间 $B_N$、诱导任务
$\bar g_N$ 和诱导核族 $\bar\cK_N$。简记
\[
c_N^\tau:=c_{\bar\cK_N}^{\tau}(\bar g_N).
\]

称系统族具有\emph{任务相对低复杂度近似闭合性质}（LC），如果存在
$\tau_N\to0$ 使
\begin{equation}
c_N^{\tau_N}=o(|B_N|).
\label{eq:LC}
\end{equation}
这要求同一列任务保真表示同时承担两件事：其状态数相对于可辨微观基线
趋零，其最优一步闭合误差也趋零。

\subsection{联合指标是前沿到原点的距离}

沿用式~\eqref{eq:diameter-frontier} 的直径前沿，定义
\begin{align}
\Gamma_N
&:=
\min_{|O_N|\le m\le|B_N|}
\max\left\{
\frac{m}{|B_N|},
\mathcal D_{\bar\cK_N,\bar g_N}(m)
\right\}
\nonumber\\
&=
\min_{\pi\,\text{细化}\,\ker\bar g_N}
\max\left\{
\frac{|\pi|}{|B_N|},
D_{\bar\cK_N}(\pi)
\right\}.
\label{eq:Gamma}
\end{align}
第一行把 $\Gamma_N$ 表示为归一化复杂度--直径前沿到原点的
$\ell^\infty$ 距离；第二行直接在全部任务保真分区上优化。两式相等：
任一大小为 $k$ 的分区在第一行取 $m=k$ 给出相同或更小的候选值；反之，
达到 $\mathcal D(m)$ 的分区实际大小不超过 $m$，在第二行给出相同或更小的
候选值。

为比较含未知宏观核的真实闭合目标，记
\begin{equation}
\Lambda_N
:=
\min_{\pi\,\text{细化}\,\ker\bar g_N}
\max\left\{
\frac{|\pi|}{|B_N|},
\eta_\pi(\bar\cK_N)
\right\}.
\label{eq:Lambda}
\end{equation}
有限分区集保证 $\Gamma_N$ 与 $\Lambda_N$ 的外层最小值都实际达到；对每个
固定候选和每个固定 $\bar K$，最优宏观核也实际达到。近似最优分区
本身不必唯一。

\begin{theorem}[压缩--闭合联合指标的表征定理]
\label{thm:gamma-characterization}
对每个 $N$ 都有定量夹逼
\begin{equation}
\frac12\Gamma_N\le\Lambda_N\le\Gamma_N.
\label{eq:gamma-lambda-sandwich}
\end{equation}
因此下列三项等价：
\begin{enumerate}
  \item 系统族满足 LC，即存在 $\tau_N\to0$ 使
  $c_N^{\tau_N}=o(|B_N|)$；
  \item 存在任务保真满射 $s_N:B_N\twoheadrightarrow Y_N$，使
  \[
  \frac{|Y_N|}{|B_N|}\to0,
  \qquad
  \sup_{\bar K\in\bar\cK_N}
  \min_{L\in\Stoch(Y_N,Y_N)}
  d_\infty(\bar KQ_{s_N},Q_{s_N}L)\to0;
  \]
  \item $\Gamma_N\to0$。
\end{enumerate}
\end{theorem}

\begin{proof}
对任意任务保真分区，令
$a_\pi:=|\pi|/|B_N|$。由引理~\ref{lem:radius-diameter}，
\[
\frac12\max\{a_\pi,D_{\bar\cK_N}(\pi)\}
\le
\max\{a_\pi,\eta_\pi(\bar\cK_N)\}
\le
\max\{a_\pi,D_{\bar\cK_N}(\pi)\}.
\]
对全部分区取最小值得式~\eqref{eq:gamma-lambda-sandwich}。

若 LC 成立，选择达到 $c_N^{\tau_N}$ 的分区，便有
$\Lambda_N\le\max\{c_N^{\tau_N}/|B_N|,\tau_N\}\to0$。反之，选择达到
$\Lambda_N$ 的分区 $\pi_N$；若 $\Lambda_N\to0$，则
$|\pi_N|/|B_N|\to0$ 且
$\eta_{\pi_N}(\bar\cK_N)\to0$。取
$\tau_N:=\eta_{\pi_N}(\bar\cK_N)$，则
$c_N^{\tau_N}\le|\pi_N|=o(|B_N|)$，即得 LC。因此 LC 等价于
$\Lambda_N\to0$，再由式~\eqref{eq:gamma-lambda-sandwich} 等价于条件 3。

最后，任务保真满射与其纤维分区一一对应，且条件 2 中的第二个量恰为该
分区的 $\eta$；所以条件 2 也等价于 $\Lambda_N\to0$。
\end{proof}

定理~\ref{thm:gamma-characterization} 的定位是表征定理：三项等价把 LC 的
存在性定义、达到的复杂度--闭合误差前沿和一列任务保真分区对应起来；其中的
定量内容是式~\eqref{eq:gamma-lambda-sandwich}，它用不含未知宏观核的纤维
预测直径在常数因子内表征真实闭合目标。该表征仍以任务保真分区为候选类；
第~\ref{sec:verification}节给出有限枚举程序及其 Bell 数规模，第~\ref{sec:construction}节
给出可绕开完整枚举的模型结构证书。

条件 2 中的最优宏观核允许随 $N$ 和 $\bar K$ 改变；整个允许
动力学族共享同一个 $L$ 对应更强的 $\min_L\sup_{\bar K}$ 问题。定理把存在性化为同一列任务保真分区上的联合
判据。若取达到 $\Gamma_N$ 的分区，则其相对状态数不超过 $\Gamma_N$，且
由式~\eqref{eq:radius-diameter}，其真实最优闭合缺陷也不超过
$\Gamma_N$；因此每个有限规模上的证书是可达到的。

\begin{corollary}[任务输出规模的必要条件]
\label{cor:attainment-output}
若系统族满足 LC，则
\[
\frac{|O_N|}{|B_N|}\to0.
\]
充分性还要求纤维预测直径同时趋零。
\end{corollary}

\begin{proof}
任务保真要求每个候选分区至少有 $|O_N|$ 个块，故
$|O_N|\le c_N^{\tau_N}=o(|B_N|)$。
\end{proof}

有限达到确立每个规模上的最优对象。计算效率、描述长度、参数量和求值成本属于独立问题。$\Gamma_N$ 的实际检验范围和估计
误差见第~\ref{sec:verification}节；附录~\ref{app:examples-counterexamples}给出 $\Gamma_N=1$ 的循环链边界例。
